% \documentclass[../Main.tex]{subfiles}
% % \documentclass{article}

% \begin{document}
\section{Motivation}
\label{section:1.1}

Secure communication has evolved through a sequence of paradigms, each introduced
to repair an operational or trust limitation left open by its predecessor. The
journey begins with symmetric cryptography and its key-distribution bottleneck,
which public-key cryptography removes at the cost of a heavy certificate
infrastructure; identity-based encryption then collapses that infrastructure but
concentrates all trust in a single key generator; threshold cryptography disperses
that trust across a committee; and the rise of blockchains adds a final
requirement, namely that confidentiality must hold only \emph{until} an ordering
has been fixed, giving rise to batched decryption. This section traces that
progression and explains why combining these ideas, in a manner that also resists
quantum adversaries, motivates the threshold batch identity-based encryption scheme
developed in this thesis.


\subsection{From Classical Cryptography to Public-Key Encryption}
\label{subsec:pke}

Modern cryptography is often traced back to Claude Shannon’s 1949 work \emph{Communication Theory of Secrecy Systems}, which established a mathematical foundation for secrecy, including the notion of perfect secrecy and the information-theoretic analysis of cryptographic systems \cite{J:Shannon49}. Shannon’s framework, however, describes \emph{symmetric} (secret-key) cryptography, in which the communicating parties must agree on a shared secret key before any message can be exchanged. While such schemes can be extremely efficient and even information-theoretically optimal, they inherit a structural obstacle known as the \emph{key-distribution problem}: a network of $n$ mutually distrusting users requires on the order of $n^2$ pairwise keys, and every one of those keys must first be transported over an already-secure channel. In open, large-scale networks such as the Internet, bootstrapping these shared secrets out-of-band is impractical, and it is precisely this obstacle that public-key cryptography was designed to overcome.

Building on this foundation, cryptography evolved from secret-key techniques into public-key cryptography (Figure~\ref{fig:PKE}), where Diffie and Hellman introduced a new paradigm for secure communication over insecure channels without requiring a pre-shared secret \cite{J:DiffieHellman76}. Their central insight was to decouple the capability to encrypt from the capability to decrypt by assigning each user a pair of mathematically related keys: a public key that may be openly published, and a private key that is never disclosed, such that knowledge of the former does not reveal the latter. Shortly afterwards this paradigm was realised by concrete constructions whose security rests on the conjectured hardness of number-theoretic problems such as integer factorization and the discrete logarithm, yielding the first practical \emph{trapdoor one-way functions} for both encryption and digital signatures. As illustrated in Figure~\ref{fig:PKE}, this traditional Public Key Encryption (PKE) paradigm relies heavily on a trusted third party, known as a Central Authority (CA) or Certificate Authority, to establish trust between communicating parties. In this model, a user (e.g., Alice) generates a cryptographic key pair consisting of a private key, which she keeps secret, and a public key. To authenticate her public key to the rest of the network, Alice groups her public key with her unique identifier (ID) and submits them to the CA through a certificate signing request. Upon verifying Alice's identity, the CA issues an X.509 Certificate. This certificate acts as a digital passport, cryptographically binding Alice's identifier to her public key, and its authenticity can in turn be traced through a chain of intermediate authorities up to a small number of widely trusted root CAs. When Alice wishes to initiate secure communication with another user, Bob, she transmits this X.509 Certificate to him. However, possessing the certificate is not sufficient to guarantee secure communication, as Alice's private key could have been compromised after the certificate was issued. To ensure the certificate is still active, Bob must check its status by querying a Validation Authority (VA). The VA stays synchronized with the CA through periodically updating messages—typically Certificate Revocation Lists (CRLs), or online status responses served through protocols such as OCSP—which contain the statuses of all issued certificates. The VA processes Bob's query and responds with a definitive "Valid or Revoked" status. Only upon receiving a "Valid" confirmation can Bob safely extract Alice's public key to encrypt messages for her or verify her digital signatures.


\begin{figure}[!htbp]
    \centering
    \includegraphics[width=1\linewidth]{Figure/Chap_1/PKE.png}
    \caption{Public Key Encryption paradigm}
    \label{fig:PKE}
\end{figure}


However, while PKE revolutionized secure communication, it introduced significant operational overhead in the form of public-key management, an apparatus collectively referred to as a Public Key Infrastructure (PKI). In traditional PKE infrastructures, the authenticity of a public key must be validated against a certificate issued by a CA, which means that trust does not flow directly between communicating users but is instead mediated by a hierarchy of authorities. This creates a dependency on complex certificate management, including the registration and distribution of public keys, the issuance and renewal of certificates, the maintenance and timely propagation of CRLs, and the storage and verification overhead of certificate chains. Each of these components is a potential weak point: a compromised or misbehaving CA can issue fraudulent certificates for arbitrary identities, revocation information may propagate slowly enough to leave a window of vulnerability, and the validation services themselves become availability-critical infrastructure. Consequently, users are burdened with the non-trivial task of verifying certificates before communication can proceed, which can introduce latency and additional points of failure into otherwise secure systems.


\subsection{Identity-Based Encryption}
\label{subsec:ibe}

To address these limitations, Shamir proposed identity-based cryptography (IBE), where a user’s public identity can serve directly as a public key, reducing the burden of certificate management~\cite{C:Shamir84}. Shamir’s original formulation posed this as a vision rather than a complete solution: he gave an identity-based signature scheme but left identity-based \emph{encryption} as an open problem, one that would remain unsolved for nearly two decades. The appeal of the idea is that it eliminates the indirection of certificates entirely. In environments where a centralized authority, e.g., a trusted dealer, is already present, the reliance on external certificate management becomes redundant. In particular, the core idea of IBE is that each joined party will receive secret keys bound to provided identifiers, such as a name, address, email, or blockchain wallet address. These keys are generated by a Private Key Generator (PKG) using a master public and secret key pair: the PKG publishes the master public key once, and on request derives, for any identity string, the corresponding identity-bound secret key through a key-extraction procedure. This approach streamlines key distribution by eliminating the need for certificates altogether; senders no longer need to retrieve and verify a certificate before encrypting a message. Instead, to encrypt messages so that the recipients can decrypt them correctly, senders simply provide the auxiliary information bound to the receivers, which serves as the users’ public key (Figure~\ref{fig:IBEprotocol}). Notably, a sender can even encrypt to an identity \emph{before} the corresponding secret key has been extracted, a property that proves essential in the batched, time-released setting discussed later in this chapter.

\begin{figure}[!htbp]
    \centering
    \includegraphics[width=1\linewidth]{Figure/Chap_1/IBEprotocol.png}
    \caption{The basic IBE protocol}
    \label{fig:IBEprotocol}
\end{figure}

Boneh and Franklin then gave the first practical identity-based encryption construction from bilinear pairings on elliptic curves, finally resolving Shamir’s open problem and turning IBE into a usable primitive for secure communication \cite{C:BonehFranklin01}. Their scheme demonstrated that the appealing certificate-free vision could be instantiated efficiently and with a rigorous security proof, and it triggered an extensive line of follow-up work on more expressive variants and stronger security models. The convenience of IBE, however, comes at a price that is the dual of the trust problem it solves. Because the PKG derives every user’s secret key from a single master secret, it is necessarily able to reproduce those keys itself, creating a \emph{key-escrow problem}: the authority that generates user private keys may also impersonate any user or decrypt any ciphertext addressed to them. The PKG therefore becomes a single, fully trusted point whose compromise is catastrophic for the entire system—precisely the kind of concentrated trust that distributed cryptography seeks to avoid.


\subsection{Threshold Cryptography and Distributed Trust}
\label{subsec:threshold}

Threshold cryptography addresses this weakness by distributing secret authority among multiple parties, following the principle of secret sharing introduced by Shamir, where a secret can be split into shares such that any sufficiently large subset can reconstruct it, while any smaller subset learns nothing about it \cite{M:Nakamoto08}. In a $(k, N)$ threshold scheme, the secret is encoded so that any $k$ of the $N$ shareholders can recover it but any $k-1$ of them obtain no information whatsoever, a guarantee that holds even against a computationally unbounded adversary. Applied to public-key primitives, this idea allows a private key to exist only \emph{implicitly}, as a quantity that is never reconstructed in one place: the parties instead compute the cryptographic operation, including decryption or signing, jointly, each contributing a partial result derived from its own share. Transplanting this mechanism onto the PKG of an identity-based scheme directly dissolves the key-escrow problem, since no single authority ever holds the master secret or any user’s full private key, and an adversary must corrupt an entire quorum rather than a single server to subvert the system. Thus, threshold identity-based encryption combines the usability of IBE with the robustness of distributed trust.

In standard encryption models, relying on a single, whole private key for a user creates a critical vulnerability: a single point of compromise exposes all their encrypted data. To mitigate this risk, a threshold key distribution scheme mathematically fragments the user's private key, distributing these shares across several independent parties. Under this architecture, an attacker attempting to disclose the user's private key to decrypt a targeted ciphertext cannot succeed by breaching a single location. Instead, they are forced to simultaneously compromise a predefined quorum of these independent parties to gather enough fragmented keys. This distributed approach dramatically elevates the difficulty of an attack, ensuring that the ciphertext remains entirely secure unless the specific threshold of fragments is successfully captured and reconstructed; conversely, it also improves availability, since the system continues to function as long as any qualified quorum of honest parties remains reachable, tolerating the failure or unavailability of a minority of authorities. The full construction is described in Figure~\ref{fig:thresholdIBE}. The operational flow begins with a Private Key Generator (PKG) that holds both the master public key $mpk$ and the master secret key $msk$. Rather than issuing full private keys directly, the PKG distributes secret key shares ($sk_1, sk_2, \dots, sk_N$) alongside the $mpk$ to $N$ independent parties ($P_1, \dots, P_N$). When a user requires their private key to decrypt a message, they broadcast their identity ($id$) to the network. A predefined threshold of honest parties (depicted as $P_1$ through $P_k$) process this request, generating partial identity-based secret keys ($sbk_1, \dots, sbk_k$). These partial keys are subsequently aggregated via a combine function to successfully reconstruct the pre-decryption key $sbk$. On the sender's side, a plaintext message $m$ is encrypted using the global $mpk$ and the recipient's $id$, producing the ciphertext $C_{id} = \text{Enc}(mpk, id, m)$. Finally, the recipient feeds the reconstructed $sbk$, the $mpk$, and the ciphertext into the decryption algorithm, $\text{Dec}(mpk, sbk, C_{id})$, to securely recover the original message $m$. Crucially, the most desirable instantiations of this flow are \emph{non-interactive}: each party derives its partial key locally from its own share and the requested identity, without any communication round among the authorities, so that the combine step is a simple public aggregation of independently produced shares.

\begin{figure}[h]
    \centering
    \includegraphics[width=1.1\linewidth]{Figure/Chap_1/thresholdIBE.png}
    \caption{The threshold IBE protocol}
    \label{fig:thresholdIBE}
\end{figure}


\subsection{Batched Decryption and Maximal Extractable Value}
\label{subsec:batch}

The advent of blockchain technology has brought a breakthrough in transparency, decentralization, and programmable trust, enabling a wide range of applications such as decentralized finance, digital identity, and autonomous governance. In these systems, users broadcast transactions that wait in a public holding area, the \emph{mempool}, until a block proposer selects, orders, and commits them into the next block. The same transparency that makes blockchain systems publicly verifiable, however, also introduces a new cryptographic challenge: Maximal Extractable Value (MEV). Because pending transactions are visible in the public mempool before execution, and because the party building a block enjoys near-total discretion over which transactions to include and in what order, that party can extract value by strategically positioning its own transactions relative to those of ordinary users. Concretely, this ordering power enables front-running (placing a transaction just ahead of a victim’s), back-running (placing one just behind it), sandwich attacks (surrounding a victim’s trade on both sides), and selective inclusion or censorship of transactions~\cite{SP:Daian20}. This issue is particularly harmful in decentralized finance, where the outcome and value of a transaction may depend heavily on its position within a block, so that the same trade can yield a materially worse price purely because of how it was ordered (Figure~\ref{fig:MEV}). 


\begin{figure}[!h]
    \centering
    \includegraphics[width=1.0\linewidth]{Figure/Chap_1/MEV.png}
    \caption{Illustration of MEV attack}
    \label{fig:MEV}
\end{figure}

Beyond the direct financial harm to users, the competition to capture MEV exerts destabilizing pressure on the consensus layer itself and encourages centralization of block production. Existing mitigations, such as private relays and private mempools, can reduce public transaction exposure, but they often shift trust to specialized intermediaries instead of eliminating the underlying ordering-information asymmetry. Therefore, a more principled cryptographic solution is needed: transaction contents should remain hidden until their ordering has already been fixed.

This motivates the rise of threshold batch identity-based encryption as a direct defense against MEV. In this approach, users encrypt their transactions before broadcasting them, while the capability to decrypt is distributed among a committee of validators or key authorities. No single party can decrypt early; only after a transaction batch has been ordered or finalized can a threshold of participants jointly release the decryption material, so that block builders are forced to commit to an ordering while still blind to the contents they are ordering. The identity-based structure makes this mechanism especially natural for blockchain settings: the encryption identity can encode a block number, epoch, application, auction round, or transaction batch label, so that ``the right to decrypt this slot'' is exactly the identity-bound key the committee releases for that slot. The \emph{batch} requirement reflects a practical efficiency concern: a block contains many transactions, and performing an independent threshold decryption for each one would be prohibitively expensive, so the goal is to release decryption for an entire batch at amortised cost. The first concrete construction in this line was proposed in \cite{EPRINT:BonLauTas25}, which combines a KZG polynomial commitment scheme with an efficient special-purpose witness encryption scheme. Its security model is comparatively weak, however, since it is proven in the generic group model. Several concurrent works have since attempted to overcome the remaining limitations and produce more concrete and efficient schemes. Boneh et al. propose an efficient BTE construction based on partial-fraction techniques, showing that algebraic batching can reduce decryption overhead while supporting short-ciphertext variants and practical implementations \cite{EPRINT:BNRT26}. Building on this direction, Agarwal et al. introduce BTX, a simpler and concretely efficient BTE scheme that is both epochless and collision-free and avoids user-chosen ciphertext indices~\cite{EPRINT:ADGRS26}.


\subsection{Toward Post-Quantum Threshold Batch IBE}
\label{subsec:pq-tbibe}

Unfortunately, despite being conceptually simple and practically efficient, all the aforementioned schemes rest on the elliptic-curve pairing assumption, whose underlying discrete-logarithm hardness is devastatingly broken by Shor’s algorithm, which solves both integer factorization and discrete logarithms in polynomial time on a quantum computer \cite{J:Shor97}. This threat is not merely asymptotic: an adversary can harvest encrypted traffic today and decrypt it once a sufficiently capable quantum machine becomes available, which is especially damaging for blockchains, where ciphertexts and commitments are recorded permanently and publicly. The cryptographic community has responded with a sustained effort to migrate to assumptions believed to resist quantum attack, and among the candidate families, lattice-based schemes have gained the broadest popularity due to their well-established post-quantum security guarantees, their reductions from worst-case lattice problems, and their versatility in supporting advanced functionalities. The security of these schemes typically reduces to the Learning With Errors (LWE) and Short Integer Solution (SIS) problems, and their most powerful constructions are enabled by \emph{lattice trapdoors}: short secret bases that permit sampling short preimages of a public matrix, the very mechanism underlying lattice-based identity-based encryption and signatures.

From this motivation, Boneh et al. \cite{EPRINT:BonLauTas25} constructed a new lattice-based batched decryption scheme, utilizing an explainable sampler and a shifted multi-preimage trapdoor sampler. More recently, Xiong, Zhang, and Chen proposed the first batched IBE from lattices in the standard model, proving security under the succinct LWE assumption \cite{EPRINT:XioChe25}. Yet making such schemes \emph{threshold} runs into a fundamental obstacle: finding an effective way to secret share the lattice trapdoor, which is the fundamental building block of essentially all lattice-based primitives, remains a non-trivial task. The difficulty is structural—a lattice trapdoor is a short basis whose usefulness depends delicately on its smallness, so naively splitting or summing trapdoor material tends either to destroy the shortness that makes it a trapdoor or to leak information about it. Prior approaches to thresholdizing lattice-based trapdoor primitives either rely on generic tools, such as the universal thresholdizer based on threshold FHE~\cite{BGG+18,ASY22}, or on MPC-based trapdoor-sharing techniques~\cite{BKP13}; however, these methods are concretely expensive, require heavy interaction between the parties, or depend on precomputed correlated randomness, none of which fits the non-interactive, public-aggregation flow desired in the mempool setting. A separate line of work develops tailored lattice-based threshold signatures~\cite{GKS24,DKM+24,EKT24,BKL+25}, but these constructions are usually interactive and rely on the specific structure of the underlying $\Sigma$-protocol, making them difficult to repurpose for general trapdoor-based primitives such as IBE. Related notions such as half trapdoors~\cite{Wee21} and subspace trapdoors~\cite{PT24} provide only weaker or different functionalities, and do not by themselves yield a distributed preimage sampler. In contrast, partial lattice trapdoors offer a direct algebraic mechanism for non-interactively distributing lattice preimage sampling: each party holds a partial trapdoor, samples a partial preimage locally, and the partial preimages recombine into a valid short preimage by simple summation, thereby giving a plug-and-play route toward threshold versions of trapdoor-based primitives such as signatures and identity-based encryption. It is precisely this primitive that the present thesis builds upon to obtain a post-quantum, non-interactive threshold batch IBE.


% \section{Related work}


\section{Our results}
\label{section:1.2}
This thesis set out to answer a concrete and timely question: \emph{can the
batched decryption used in encrypted mempools be made both post-quantum and
trustless at the same time?} Existing batch threshold encryption schemes achieve
the trustless property but rely on pairing-based assumptions that do not survive
quantum adversaries, while the lattice machinery that would make them
post-quantum, namely trapdoor preimage sampling, has long resisted efficient,
non-interactive thresholdization. The central contribution of this thesis is a
construction that resolves this tension: a lattice-based threshold batch
identity-based encryption (\textsf{TBIBE}) scheme in which decryption authority
is split $k$-out-of-$N$ across a committee, without any interaction between the
parties at decryption time.
 
Concretely, the work of this thesis can be summarized in the following
contributions.
 
\begin{enumerate}
    \item \textbf{A new post-quantum threshold batch IBE scheme.}
    We design a $k$-out-of-$N$ threshold batch IBE by combining the index-dependent
    batched decryption scheme of Boneh, Laufer and Tas~\cite{EPRINT:BonLauTas25}
    with the partial lattice trapdoors of Albrecht, Lai, Lapiha and
    Woo~\cite{AC:ALLW25}. The key technical step is to replace the single
    $\mathsf{SamplePre}$ call that produces the batch pre-decryption key
    $\mathsf{sbk}$ with its distributed analogue $\mathsf{PSampPre}$ followed by
    reconstruction $\mathsf{Rec}$, so that each authority holds only a partial
    trapdoor $\mathsf{ptd}_j$, locally emits a partial key $\mathsf{sbk}_j$, and any
    qualified quorum recombines these shares into the full key by simple summation.
    To our knowledge, this is the first construction to thresholdize a lattice-based
    \emph{batched} IBE in this purely algebraic, non-interactive way.
 
    \item \textbf{A formal syntax and security model for $\mathsf{TBIBE}$.}
    We give a precise five-algorithm syntax
    ($\mathsf{Setup}$, $\mathsf{Enc}$, $\mathsf{PreDec}$, $\mathsf{Combine}$,
    $\mathsf{Dec}$) and a correctness definition for the threshold setting, and we
    introduce a \emph{highly selective} security notion that faithfully captures
    what the secret-sharing layer permits an adversary to learn, including the
    corruption of up to $k-1$ parties and the release of all but one partial key
    for the challenge quorum.
 
    \item \textbf{Correctness and security proofs.}
    We prove correctness from the structural properties of the shifted
    multi-preimage sampler, the explainable $\mathsf{SampleLeft}$ procedure, and the
    reconstruction guarantee of the partial trapdoor scheme. We then prove security
    through an explicit sequence of hybrid games that reduces the advantage of any
    efficient adversary to the
    $(nk, 2mk, \chi_0, \chi_1, \sigma_2)$-indistinguishability of the partial
    trapdoor, and hence to the $\kappa$-LWE assumption (with one-wayness resting on
    $\kappa$-SIS), closing the remaining gaps with Gaussian-smoothing and
    noise-flooding arguments.
 
    \item \textbf{A unified treatment and analysis of the building blocks.}
    We assemble and present, in a single consistent framework, the lattice
    tools the construction depends on, gadget and Ajtai trapdoors, the explainable
    sampler, the shifted multi-preimage trapdoor sampler, and partial lattice
    trapdoors, and we adapt their correctness and hardness statements to the
    dimensions used by our scheme.
 
    \item \textbf{An efficiency study and proof-of-concept evaluation.}
    We derive the asymptotic sizes of the public key, ciphertext, and pre-decryption
    key, analyze the runtime of each algorithm, and identify the pre-decryption phase
    (explainable $\mathsf{SampleLeft}$ and $\mathsf{PSampPre}$) as the dominant cost.
    A proof-of-concept implementation under representative parameters demonstrates
    that the scheme is feasible and isolates the concrete bottlenecks that future
    optimization should target.
\end{enumerate}
 
Taken together, these contributions show that partial lattice trapdoors are an
effective and practical tool for distributing decryption authority in
lattice-based identity-based encryption. The resulting scheme inherits the most
valuable property of its underlying trapdoor-sharing mechanism, namely that the
distribution of the decryption capability is fully \emph{non-interactive}: parties
never communicate to produce their shares, which sets the construction apart from
prior thresholdization routes based on threshold FHE or multiparty computation.




\label{section:1.4}

\section{Thesis Organization}

The remainder of this thesis is organized as follows.

\begin{itemize}
    \item \textbf{Chapter~2} presents the necessary preliminaries for the proposed construction. 
    It reviews fundamental lattice theory, lattice-based hardness assumptions, gadget and 
    Ajtai trapdoors, explainable sampler, shifted multi-preimage trapdoor 
    samplers, and partial lattice trapdoors.

    \item \textbf{Chapter~3} introduces the proposed threshold batch identity-based encryption 
    construction. It first defines the syntax of the scheme and then describes the setup, 
    encryption, partial pre-decryption, key-combination, and decryption algorithms in detail.

    \item \textbf{Chapter~4} analyzes the correctness and security of the proposed construction. 
    It proves that the decryption procedure correctly recovers the encrypted messages under 
    the chosen parameters, derives the asymptotic sizes of the main objects, and establishes 
    security in the highly selective model through a sequence of hybrid games.

    \item \textbf{Chapter~5} discusses the experimental evaluation of the proposed scheme, 
    including implementation considerations and the efficiency of the main algorithms.

    \item \textbf{Chapter~6} concludes the thesis by summarizing the main contributions and 
    discussing possible directions for future work.

    \item \textbf{Appendix} provides additional technical details, including the shifted 
    multi-preimage trapdoor sampler, the construction of $\mathsf{SampleLeft}$, the explainable 
    $\mathsf{SampleLeft}$ procedure, and the partial gadget trapdoor construction.
\end{itemize}


% \end{document}